\documentclass[12pt,a4paper]{article}

\usepackage[a4paper,margin=2.2cm]{geometry}
\usepackage{fontspec}
\usepackage{polyglossia}
\usepackage{enumitem}
\usepackage{hyperref}

\setmainlanguage{russian}
\setotherlanguage{english}

\setmainfont[
  Path=../fonts/,
  UprightFont=pt-astra-serif_regular.ttf,
  BoldFont=pt-astra-serif_bold.ttf,
  ItalicFont=pt-astra-serif_italic.ttf,
  BoldItalicFont=pt-astra-serif_bold-italic.ttf
]{PT Astra Serif}

\setmonofont[
  Path=../fonts/fira-code/,
  UprightFont=FiraCode-Regular.ttf,
  BoldFont=FiraCode-Bold.ttf
]{Fira Code}

\hypersetup{
  hidelinks,
  pdfauthor={Нахатакян Артур Романович},
  pdftitle={План демонстрации программного продукта}
}

\setlist[itemize]{leftmargin=1.4cm}
\setlist[enumerate]{leftmargin=1.4cm}
\setlength{\parindent}{0pt}
\setlength{\parskip}{0.55em}
\emergencystretch=2em

\begin{document}

\begin{center}
{\Large \textbf{План демонстрации программного продукта}}\\[0.8em]
{\normalsize Тема ВКР: «Разработка системы управления контентом в сфере маркетинга на базе Strapi и Astro»}
\end{center}

\section*{Назначение документа}

Документ фиксирует план демонстрации разработанного программного продукта без внесения
изменений в код проекта. План опирается на уже реализованный и подтвержденный baseline
системы и предназначен для предзащиты ВКР.

\section*{Цель демонстрации}

Показать, что разработанная система:
\begin{itemize}
  \item реализует \texttt{CMS-first} подход к управлению контентом;
  \item поддерживает раздельные контуры публичной витрины и административного управления;
  \item обеспечивает локализацию, предварительный просмотр, базовый \texttt{SEO}-контур и
  формы взаимодействия с пользователем;
  \item имеет воспроизводимый публикационный контур
  \texttt{publish -> webhook -> rebuild}.
\end{itemize}

\section*{Формат демонстрации}

Рекомендуемая длительность: \texttt{10--15 минут}.

Основной формат:
\begin{enumerate}
  \item краткое введение в архитектуру;
  \item показ административного контура \texttt{Strapi};
  \item показ публичной витрины \texttt{Astro};
  \item показ preview-сценария;
  \item краткий разбор эксплуатационного и публикационного контура;
  \item фиксация ограничений текущей версии.
\end{enumerate}

\section*{Среда демонстрации}

\subsection*{Требуемые сервисы}

\begin{itemize}
  \item frontend runtime: \texttt{http://localhost:4321};
  \item CMS runtime: \texttt{http://localhost:1337};
  \item при необходимости build-preview: \texttt{http://localhost:4322}.
\end{itemize}

\subsection*{Минимальная проверка перед показом}

Перед демонстрацией рекомендуется убедиться, что:
\begin{itemize}
  \item открывается главная страница \texttt{/ru/};
  \item открывается английская версия \texttt{/en/};
  \item доступна \texttt{Strapi} admin panel;
  \item работает хотя бы один preview-сценарий для \texttt{page};
  \item в CMS есть опубликованные записи для \texttt{page}, \texttt{article},
  \texttt{project}, \texttt{vacancy};
  \item в CMS есть хотя бы один draft для показа preview;
  \item заранее известны логины ролей, если планируется показ разделения доступа.
\end{itemize}

\subsection*{Что не требуется}

\begin{itemize}
  \item не требуется вносить изменения в код проекта;
  \item не требуется live-доработка данных во время выступления;
  \item не требуется запуск полного production-пути с внешним \texttt{Vercel} rebuild.
\end{itemize}

\section*{Сценарий демонстрации}

\subsection*{1. Введение и постановка задачи}

Время: \texttt{1--2 минуты}.

Что показать и проговорить:
\begin{itemize}
  \item тему ВКР и назначение системы;
  \item проблему кодо-центричного управления маркетинговым сайтом;
  \item идею разделения на \texttt{Strapi} как контентное ядро и \texttt{Astro} как
  публичную витрину.
\end{itemize}

Ожидаемый результат: комиссия понимает, зачем в проекте одновременно нужны \texttt{CMS},
публичный frontend, preview и публикационный контур.

\subsection*{2. Архитектура и контуры системы}

Время: \texttt{1--2 минуты}.

Что показать:
\begin{itemize}
  \item краткую схему взаимодействия \texttt{Strapi -> API -> Astro};
  \item два основных контура:
  \begin{itemize}
    \item маркетинговый: \texttt{global}, \texttt{home-page}, \texttt{page},
    \texttt{article}, \texttt{author}, \texttt{project}, \texttt{lead-submission};
    \item карьерный: \texttt{vacancy}, \texttt{industry}, \texttt{job-role},
    \texttt{vacancy-application}.
  \end{itemize}
\end{itemize}

Что проговорить:
\begin{itemize}
  \item роли \texttt{administrator}, \texttt{marketer/content-manager}, \texttt{editor},
  \texttt{hr};
  \item принцип, что публичная витрина не пишет в \texttt{Strapi} напрямую из браузера.
\end{itemize}

Ожидаемый результат: показано, что система спроектирована как разделенная управляемая
платформа.

\subsection*{3. Демонстрация CMS-контура}

Время: \texttt{2--3 минуты}.

Что показать в \texttt{Strapi}:
\begin{itemize}
  \item список ключевых сущностей;
  \item структуру \texttt{home-page} или \texttt{page};
  \item наличие \texttt{Dynamic Zone} у страниц;
  \item \texttt{SEO}-поля у страницы или детальной сущности;
  \item наличие опубликованного и чернового контента;
  \item при возможности --- разделение ролей или описание того, кто что редактирует.
\end{itemize}

Что проговорить:
\begin{itemize}
  \item \texttt{editor} готовит черновик, но не публикует;
  \item \texttt{marketer/content-manager} публикует маркетинговый контент;
  \item \texttt{hr} управляет вакансиями и откликами.
\end{itemize}

Ожидаемый результат: показано, что контент редактируется в \texttt{CMS}, а не в исходном
коде frontend.

\subsection*{4. Демонстрация публичной витрины}

Время: \texttt{3--4 минуты}.

Что открыть в браузере:
\begin{enumerate}
  \item \texttt{/ru/} --- главная страница;
  \item \texttt{/en/} --- английская версия главной;
  \item representative \texttt{page}, например \texttt{/ru/cms-first-platform/};
  \item список или detail статьи \texttt{/ru/articles/...};
  \item список или detail проекта \texttt{/ru/projects/...};
  \item \texttt{/vacancies/} и detail вакансии.
\end{enumerate}

Что проговорить:
\begin{itemize}
  \item storefront-core работает в локалях \texttt{ru/en};
  \item страницы, статьи и проекты приходят из \texttt{CMS};
  \item вакансии вынесены в отдельный прикладной модуль;
  \item frontend выполняет роль публичной витрины и использует данные \texttt{Strapi}.
\end{itemize}

Ожидаемый результат: комиссия видит, что система реально публикует контентные разделы, а
не ограничивается только админкой.

\subsection*{5. Демонстрация preview-сценария}

Время: \texttt{2 минуты}.

Что показать:
\begin{itemize}
  \item preview action для черновой \texttt{page} или другой сущности;
  \item переход через \texttt{/api/preview};
  \item открытие маршрута \texttt{/preview/...};
  \item отличие preview от обычной публичной страницы.
\end{itemize}

Что проговорить:
\begin{itemize}
  \item preview работает без раскрытия публичного draft API;
  \item draft-чтение выполняется через \texttt{PREVIEW\_SECRET};
  \item preview-страницы получают \texttt{noindex}.
\end{itemize}

Ожидаемый результат: показан один из ключевых инженерных результатов --- безопасный
предпросмотр черновика.

\subsection*{6. Демонстрация форм взаимодействия}

Время: \texttt{1--2 минуты}.

Безопасный вариант показа:
\begin{itemize}
  \item открыть публичную форму лида или отклика;
  \item показать обязательные поля;
  \item продемонстрировать валидацию на невалидных данных;
  \item при необходимости кратко показать, что отправка идет через server-side route, а не
  напрямую в \texttt{Strapi}.
\end{itemize}

Если нежелательно изменять данные в системе, следует ограничиться показом формы и сценария
валидации без финальной успешной отправки.

Если допустима мутация тестовых данных, можно показать успешную отправку и наличие записи в
\texttt{Strapi}.

Ожидаемый результат: показано, что прикладные формы встроены в систему и защищены
server-side слоем.

\subsection*{7. Демонстрация публикационного контура}

Время: \texttt{1--2 минуты}.

Что показать:
\begin{itemize}
  \item наличие managed webhook в \texttt{Strapi};
  \item подписку на события \texttt{entry.publish} и \texttt{entry.unpublish};
  \item краткое объяснение цепочки \texttt{publish -> webhook -> rebuild}.
\end{itemize}

Что проговорить:
\begin{itemize}
  \item локально подтвержден факт существования и синхронизации webhook;
  \item полный внешний production-путь с \texttt{Vercel} не требуется демонстрировать
  вживую;
  \item для целей ВКР важно показать сам инженерный контур и его воспроизводимость.
\end{itemize}

Ожидаемый результат: комиссия видит, что публикация контента связана с обновлением
публичной витрины.

\subsection*{8. Итог и ограничения}

Время: \texttt{1 минута}.

Что проговорить в завершении:
\begin{itemize}
  \item система реализует \texttt{CMS-first} модель для маркетингового сайта;
  \item контентное управление отделено от публичной витрины;
  \item реализованы локализация, preview, формы и публикационный контур;
  \item проект имеет воспроизводимый набор проверок и acceptance-сценариев.
\end{itemize}

Какие ограничения нужно проговаривать честно:
\begin{itemize}
  \item английское наполнение detail-разделов \texttt{articles/projects} зависит от
  фактического dataset;
  \item текущий \texttt{SEO} dataset требует дополнительной финальной проверки индексации
  публичных страниц;
  \item browser-level аудит через \texttt{Lighthouse/axe} не является обязательной частью
  текущей live-демонстрации;
  \item live-показ внешнего \texttt{Vercel} rebuild не обязателен.
\end{itemize}

\section*{Что лучше не показывать как основной сценарий}

Во время демонстрации не стоит делать центральным элементом показа:
\begin{itemize}
  \item редактирование production-данных без заранее подготовленного тестового сценария;
  \item зависимые от внешней инфраструктуры действия, которые могут сорваться из-за сети;
  \item англоязычные detail-страницы \texttt{articles/projects}, если они не заполнены в
  текущем dataset;
  \item спорные \texttt{SEO}-утверждения без предварительно проверенного фактического
  состояния.
\end{itemize}

\section*{Резервный план}

Если часть live-сценариев не открывается во время демонстрации, рекомендуется:
\begin{enumerate}
  \item сохранить основную логику показа через уже работающие маршруты \texttt{/ru/},
  \texttt{page}, \texttt{preview}, \texttt{vacancies};
  \item использовать заранее открытые вкладки браузера;
  \item опираться на уже подготовленные диаграммы, acceptance-материалы и структуру
  сущностей в \texttt{Strapi};
  \item делать акцент на архитектуре, редакторском workflow и воспроизводимых проверках.
\end{enumerate}

\section*{Вывод}

Для предзащиты оптимален сценарий, в котором демонстрация строится не вокруг внесения
новых изменений в проект, а вокруг уже реализованных и подтвержденных контуров:
контентного управления, публичной витрины, preview, форм и публикации. Это позволяет
показать зрелость решения и избежать риска live-ошибок, связанных с ненужной доработкой
проекта непосредственно перед выступлением.

\end{document}
